\cotefile{156-3}
\scanpage{26}\mspage{25}
% L1–4
pour
\[
\partial T=\text{bord absolu de }T
=\left\{\deleted{t\in T\mid t\text{ est }}\begin{aligned}
&\text{un plus petit}\\[-2pt]
&\text{ou plus grand}\\[-2pt]
&\text{élément de }T
\end{aligned}\right\}
\]
% L5–6
\deleted{Si $T$ est réduit à un pt,}

On a donc
% L7–10
\[
\begin{aligned}
\partial\{a\}&=\{a\}\\
\partial[\unsure{o|a|0},b]&=\{a,b\} &&\text{si }a\leq b\\
\partial L_{\geq a}&=\partial L_{\leq a}=\{a\} &&\forall\,a\in L\\
\deleted{\unsure{\partial L=\varnothing}}\quad\partial L&=\partial L\quad\text{(sic)}
\end{aligned}
\]
\margin{(sans exclure le cas $a\in\partial L$)}
% L11–12
Notons que
\[
T^\bullet\cap T\subset\partial T\subset\deleted{\partial L}(T^\bullet\cap T)\cup\partial L
\]
% L13–15
i.e tout $t\in T$ qui est pt frontière, est dans le bord \inserted{de $T$}, et l'inverse est vrai si $t\notin\partial L$. Plus précisément
% L16–18
\[
\begin{aligned}
&\deleted{\text{Si }\operatorname{Card}\,T=1_{\unsure{\mathrm{elt}|\mathrm{él}}},\quad\partial T=T^\bullet=T}\\
&\deleted{\text{Sinon}}\quad
\left\{\begin{aligned}
&\deleted{T^\bullet\cap T=\partial T\setminus(\partial T\cap\partial L)}\quad\deleted{\text{\unsure{i.e.}}}\\
&\deleted{\partial T=(T^\bullet\cap T)\cup T\cap\partial L}
\end{aligned}\right.
\end{aligned}
\]
% L19
\[
\boxed{\partial T=(T^\bullet\cap T)\mathbin{\underline{\cup}}(T\cap\partial L)}
\]
% L20–22
\unsure{Tronçons|Tronçon} : C'est une partie \deleted{\unsure{de la réduite|de la réunion}} \deleted{de la forme} \deleted{\unsure{convexe}} de la forme $T'=T\setminus\alpha$, avec $\alpha\subset\partial T$.

\deleted{\unsure{Tronçon ouvert|Tronçons ouverts} :} c'est un ens de la forme
% L23–26
\[
\left\{\begin{aligned}
&]a,b],\;[a,b[,\;]a,b[,\quad(a<b)\\
&L_{<a},\ \deleted{\unsure{a\in\partial L})}\quad(a\neq\omega)\\
&L_{>a}\hspace{7em}(a\neq\alpha)\\
&L
\end{aligned}\right.
\]
% slanted
% L lower left, oblique
\margin{\uline{NB} Cette écriture est \uline{unique} on pose $T=\widetilde{T'}$ (c'est une \unsure{sorte} d'adhérence \unsure{de $T$|de $T'$} -- \unsure{c'est} le \unsure{sous-tronçon|seul tronçon} le plus petit fermé \unsure{qui contient} $T'$.}
